
\section{Reproducibility}
\label{sec:reproducibility}

The production workflow uses frozen base seed \texttt{20260503} and the exact
seed namespace \texttt{ac-control-validation-m-surrogate-v1}. Randomized seeds
are deterministic unsigned integers derived from the first eight bytes of
SHA256 applied to the namespace, base seed, canonical galaxy name, tier name,
and one-indexed iteration. Python's built-in hash is not used. Circular shifts
and fixed references do not use random seeds.

The release plan follows general scientific-computing and
reproducible-computational-research guidance: record executable entry points,
inputs, outputs, environment details, and provenance sufficient for independent
audit \citep{Wilson2014,Sandve2013}.

The recorded environment includes Python \texttt{3.12.7}, NumPy
\texttt{1.26.4} \citep{Harris2020}, and code commit \texttt{c8a6c950bd9e}. The production entry
point is \texttt{python src/\allowbreak run\_m\_surrogate\_extension.py
\allowbreak --production}.
Invalid realizations were retained and counted in emitted totals while
excluded from valid empirical statistics. Detailed audit records are retained
in repository notes and supplementary provenance files.



